\chapter{Implementacja modułów funkcjonalnych}\label{ch:implementacja}

\section{Moduł zarządzania regułami blokującymi i zezwalającymi na ruch (porty i adresy IP)}\label{sec:modul-reguly}

\subsection{Interaktywny kreator reguł przychodzących (Inbound) i wychodzących (Outbound)}\label{subsec:kreator-regul}

\subsection{Walidacja danych wejściowych i tworzenie nowych reguł zapory}\label{subsec:implementacja-new-rule}

\subsection{Przeglądanie, filtrowanie i usuwanie istniejących reguł}\label{subsec:przegladanie-usuwanie-regul}

\section{Moduł automatycznego włączania i wyłączania reguł w zależności od profilu sieci}\label{sec:modul-profile}

\subsection{Detekcja aktywnego profilu sieciowego systemu}\label{subsec:detekcja-profilu}

\subsection{Dynamiczna aktywacja i dezaktywacja reguł dla profili Domain, Private i Public}\label{subsec:aktywacja-regul-profil}

\section{Moduł logowania wszystkich zmian w polityce zapory}\label{sec:modul-logowanie}

\subsection{Rejestrowanie operacji dodawania, modyfikacji i usuwania reguł}\label{subsec:rejestrowanie-operacji}

\subsection{Struktura pliku dziennika audytu i identyfikacja kontekstu wykonania (kto, co, kiedy)}\label{subsec:struktura-pliku-audytu}

\section{Moduł tworzenia kopii zapasowej i przywracania ustawień zapory}\label{sec:modul-backup}

\subsection{Eksport bieżącej polityki i reguł zapory do pliku kopii zapasowej}\label{subsec:eksport-konfiguracji}

\subsection{Procedura przywracania stanu zapory i weryfikacja spójności kopii}\label{subsec:import-konfiguracji}
